\documentclass[journal]{IEEEtran}
\usepackage[utf8]{inputenc}
\usepackage{amsmath, amssymb, amsfonts}
\usepackage{graphicx}
\usepackage{booktabs}
\usepackage{cite}
\usepackage{subcaption}

\title{Robust Neural-Algorithmic Framework for Directed $t$-Spanner Construction in Metropolitan Road Networks: Integrating Meta-Evolutionary Optimization, Epistemic Uncertainty, and Continual Learning}
\author{Soheyl Falahzade}

\begin{document}
\maketitle

\begin{abstract}
Real-time routing in Intelligent Transportation Systems (ITS) is severely bottlenecked by the $O(m \cdot (n + m \log n))$ computational complexity associated with constructing exact geometric $t$-spanners on large-scale urban topologies. While machine learning has emerged as a promising paradigm to bypass these computational limits, standard neural predictors lack formal topological guarantees and rely on a single hard decision threshold that ignores the structural importance of individual edges. This paper introduces a robust hybrid neural-algorithmic framework that (i) shifts the heavy computational burden of spanner construction to an offline representation learning phase while enforcing strict structural constraints online via a cutoff-bounded directed repair procedure, and (ii) replaces the conventional hard-threshold pruning decision with a Mamdani-style fuzzy inference system that jointly reasons over the network's epistemic uncertainty (via Monte Carlo Dropout) and each edge's structural importance (via Edge Betweenness Centrality). Experimental evaluations across three heterogeneous urban fabrics (Eindhoven, Paris, and Rome), each validated over 15--18 independent stochastic seeds at matched sparsification levels, show that the proposed fuzzy pruning mechanism consistently requires fewer post-hoc geometric repairs than both a genetic-algorithm-tuned deterministic threshold with active self-correcting feedback (7.1--16.6\% fewer repairs) and a sparsity-matched random pruning baseline (0.4--10.2\% fewer repairs), across every tested city without exception. We further report, with full transparency, that a single-parameter threshold search does not by itself benefit from evolutionary optimization over random search under equal evaluation budget -- underscoring that the genuine value of evolutionary search in this framework lies in the higher-dimensional joint optimization of loss weights and calibration parameters, not in scalar threshold tuning. Under extreme network corruption, the uncertainty-aware framework maintains full Strong Connectivity (SCC) and a bounded directed stretch factor across all tested configurations.
\end{abstract}

\begin{IEEEkeywords}
Intelligent Transportation Systems, Graph Neural Networks, $t$-Spanner, Strongly Connected Components, Epistemic Uncertainty, Continual Learning, Genetic Algorithms, High-Performance Computing.
\end{IEEEkeywords}

\section{Introduction}
Modern navigation engines, dynamic fleet management systems, and cooperative autonomous vehicles rely heavily on spatial graph sparsification to minimize the search space of $A^*$ or Dijkstra-based routing agents. In the domain of computational geometry, a directed $t$-spanner $H = (V, E')$ of a spatial graph $G = (V, E)$ is a sparse subgraph that preserves the shortest path distances between any pair of vertices within a strict multiplicative stretch factor $t$:
\begin{equation}
    d_H(u, v) \le t \cdot d_G(u, v) \quad \forall u, v \in V
\end{equation}

In intelligent transportation systems (ITS), maintaining the topological integrity of directed networks is extremely challenging due to asymmetric edge weights arising from one-way streets, turn restrictions, and dynamic traffic congestion. Traditional routing engines struggle with dynamic overheads. Recent paradigm shifts in intelligent transportation systems emphasize the integration of deep learning estimators for real-time route planning on asymmetric topologies [18], yet enforcing strict geometric bounds remains an unresolved challenge.

Historically, constructing high-quality geometric spanners has relied on the classic Greedy Spanner algorithm, which exhibits a computational complexity of $O(m \cdot (n + m \log n))$ using standard Dijkstra implementations. While theoretically optimal in terms of degree, weight, and edge count, this iterative construction is computationally prohibitive for metropolitan-scale directed networks. For instance, in a dense city map with over 40,000 nodes, the classic greedy algorithm requires hours of computation, making real-time map updates or online route optimization impossible.

To address this challenge, researchers have recently turned to Graph Neural Networks (GNNs) to predict edge importance and sparsify road networks in a data-driven manner. However, pure deep learning methods suffer from a critical flaw: they lack mathematical guarantees. A GNN may classify a critical highway bridge as ``prunable'' based on its local neighborhood features, resulting in infinite stretch (disconnection) or catastrophic detours. This risk is further amplified when the neural model is exposed to out-of-distribution (OOD) city topologies or noisy spatial data.

In this paper, we propose a robust, hybrid neural-algorithmic framework, named \textit{Geo-SmartSpanner}, which bridges the gap between statistical learning and hard geometric constraints. Our framework uses a scale-invariant GNN with Monte Carlo Dropout to propose candidate pruning lists with well-calibrated confidence scores. To guarantee the $t$-spanner property, we implement a parallelized, directed, batch-sequential geometric repair module that intercepts GNN decisions only when a theoretical violation is detected. Furthermore, we address the challenge of topological domain adaptation across different cities using a continual learning memory buffer.

The key contributions of this paper are summarized as follows:
\begin{itemize}
    \item We present a decoupled neural-algorithmic framework that shifts the complexity of directed $t$-spanner construction from $O(m \cdot (n + m \log n))$ to $O(n)$ offline learning and sub-second online inference.
    \item We introduce an active geometric feedback loop that utilizes the failures of the neural network (repaired edges) to fine-tune the GNN parameters, creating a self-correcting geometric model.
    \item We propose a directed Bounded-Stretch GLV-Repair algorithm that maintains 100\% Strongly Connected Components (SCC) and strictly bounds the directed stretch factor below 1.5 under extreme spatial noise.
    \item We integrate a meta-evolutionary optimization layer based on Genetic Algorithms (GA) to systematically search for Pareto-optimal decision thresholds, eliminating empirical trial-and-error.
\end{itemize}

---

\section{Related Work}
The study of geometric networks and spanner construction was pioneered by several seminal works. Graph spanners were first introduced by Peleg and Schaffer [1] in the context of distributed computing, and extended to planar structures by Chew [2]. Smid [3] provided a comprehensive analysis of closest-point problems and spanner networks. Farshi and Gudmundsson [6] provided the first extensive experimental evaluation of geometric $t$-spanners, highlighting the trade-offs between different quality measures such as degree, weight, and diameter. Bose et al. [7] subsequently reduced the naive cubic time complexity of the greedy spanner to near-quadratic bounds. However, as demonstrated by Bakhshesh and Farshi [8], even optimized variants like the gap-greedy spanner suffer from substantial computational overhead in practice, requiring several minutes to process small point sets.

The theoretical complexity of spanner design in directed graphs (oriented spanners) was recently explored by Buchin et al. [13], who proved that computing oriented spanners with optimal dilation is NP-Hard on general directed graphs. To address this computational bottleneck in practice, recent literature has shifted towards heuristic approximation. In 2026, Buchin, Rehs, and Scheele [14] introduced the concept of constructing ``doppelgangers'' of greedy geometric spanners, seeking fast algorithms that can mimic the sparseness of the greedy spanner in polynomial time. Recent advancements have expanded fault-tolerant spanner designs to planar structures [12] and proposed greedy strategies for region-fault resilient geometric networks [11], yet their computational cost under real-time directed constraints remains a significant bottleneck.

While machine learning has recently been applied to network sparsification, existing neural graph models are typically transductive and fail to guarantee connectivity. Our work differs from these previous approaches by combining the inductive generalization of Graph Neural Networks with the exact, directed verification of computational geometry, providing a highly scalable and mathematically robust solution to the directed $t$-spanner problem.

---

\section{Proposed Methodology}

Our framework consists of three tightly coupled components designed to ensure topological safety and high-performance inference under directed traffic constraints.

\subsection{Uncertainty-Aware Neural Edge Classification}
To classify edges, we implement an inductive \texttt{GeometricEdgeSAGE} architecture, building upon the inductive framework proposed by Hamilton et al. [20]. To ensure generalization across cities with different coordinate systems, we engineer scale- and symmetry-invariant spatial features for each edge $i = (u, v) \in E$:
\begin{equation}
    f_i = \left[ \bar{d}(u,v), deg_{in}(u), deg_{out}(v), \mathcal{C}(u, v) \right]
\end{equation}
where $\bar{d}(u, v)$ is the locally normalized road segment length, $deg$ represents node degrees, and $\mathcal{C}$ denotes the local edge betweenness centrality.

Instead of relying on deterministic predictions, which are vulnerable to noise, we estimate the epistemic uncertainty of our GNN using Monte Carlo (MC) Dropout. During online inference, we perform $M$ stochastic forward passes by keeping the dropout layers active:
\begin{equation}
    \mu(i) = \frac{1}{M} \sum_{j=1}^{M} f_\theta(x_i)_j
\end{equation}
\begin{equation}
    \sigma^2(i) = \frac{1}{M} \sum_{j=1}^{M} \left( f_\theta(x_i)_j - \mu(i) \right)^2
\end{equation}

The pruning decision is then governed by the risk-averse calibrated probability:
\begin{equation}
    \hat{y}_i = \mu(i) + \lambda \cdot \sigma(i)
\end{equation}
An edge $i$ is added to the candidate pruning list $E_{prune}$ if and only if $\hat{y}_i \le \tau$, where $\tau = 0.40$ represents the calibrated decision threshold.

\subsection{Topological Safeguards \& Directed GLV-Repair}
To prevent catastrophic structural disconnections, we first apply a hard topological lock. Let $E_{lock}$ be the set of edges incident to low-degree boundary vertices:
\begin{equation}
    E_{lock} = \{ (u, v) \in E \mid deg_{in}(v) \le 1 \lor deg_{out}(u) \le 1 \}
\end{equation}
These edges are strictly protected and never pruned, preserving the basic connectivity of peripheral neighborhoods.

The remaining candidate edges in $E_{prune} \setminus E_{lock}$ are processed using a \textit{Batch-Sequential GLV-Repair} algorithm. In undirected graphs, pruning decisions can be made independently. However, in directed graphs, pruning can cause asymmetric routing bottlenecks, a phenomenon we refer to as the \textit{Blind Batch-Pruning Paradox}.

To resolve this paradox, we sort $E_{prune}$ by edge length and divide it into batches of size $B = 250$. For each batch, we evaluate the shortest path on the remaining sparse graph $H$ using parallelized, directed SciPy Dijkstra. For each pruned edge $e = (u, v) \in E_{prune}$, we check:
\begin{equation}
    d_H(u, v) \le t \cdot w(e)
\end{equation}
where $t = 1.5$ is the target stretch limit. If this condition is violated, $e$ is restored to $H$. Because we process in batches, any restored edge in batch $k$ immediately serves as a valid detour for subsequent batches, preventing redundant edge restorations.

\subsection{Meta-Evolutionary Threshold Selection}
To avoid arbitrary thresholding and systematic biases, we deploy a Genetic Algorithm (GA) to find the Pareto-optimal values for $\lambda$, $\tau$, and loss weights. While methods like Particle Swarm Optimization (PSO) and Differential Evolution (DE) are suitable for continuous fields, the highly discontinuous, discrete nature of topological graph mutations makes them prone to premature convergence. The crossover and mutation operators of the GA allow it to maintain population diversity, successfully navigating the non-convex multi-objective landscape defined by:
\begin{equation}
    \text{Fitness} = w_1 \cdot \text{Sparsification} - w_2 \cdot \max_{u,v} \left( \frac{d_H(u,v)}{d_G(u,v)} - t \right)
\end{equation}
\subsection{Fuzzy-Inference Pruning Decision}
\label{sec:fuzzy}
 
A hard decision threshold on the calibrated pruning probability $\hat{y}_i$ (Eq. 5) treats all edges with $\hat{y}_i$ near the boundary identically, regardless of their structural role in the network. To address this, we replace the hard threshold with a Mamdani-style fuzzy inference system that combines two normalized inputs per edge $i$: the calibrated probability $\hat{y}_i \in [0,1]$ and the edge betweenness centrality $C(u,v) \in [0,1]$ (min-max normalized per city).
 
Each input is partitioned into three linguistic terms (\textit{Low}, \textit{Medium}, \textit{High}) using triangular membership functions:
\begin{equation}
    \mu_{\text{Low}}(x) = \max\left(0, \min\left(\frac{x - a_1}{b_1 - a_1}, \frac{c_1 - x}{c_1 - b_1}\right)\right)
\end{equation}
with analogous definitions for $\mu_{\text{Medium}}$ and $\mu_{\text{High}}$, parameterized by knot triples $(a,b,c)$ fitted such that the three terms overlap smoothly across $[0,1]$.
 
We define a rule base of five Mamdani rules that encode the intuition that an edge should be retained either when the model is confident it is necessary, or when the model is uncertain but the edge occupies a structurally critical position:
\begin{align}
    R_1&: \text{IF } \hat{y}_i \text{ is High} &\Rightarrow \text{KEEP} \\
    R_2&: \text{IF } \hat{y}_i \text{ is Medium AND } C(u,v) \text{ is High} &\Rightarrow \text{KEEP} \\
    R_3&: \text{IF } \hat{y}_i \text{ is Medium AND } C(u,v) \text{ is Low} &\Rightarrow \text{PRUNE} \\
    R_4&: \text{IF } \hat{y}_i \text{ is Low AND } C(u,v) \text{ is High} &\Rightarrow \text{KEEP (cautious)} \\
    R_5&: \text{IF } \hat{y}_i \text{ is Low AND } C(u,v) \text{ is Low} &\Rightarrow \text{PRUNE}
\end{align}
 
Each rule's firing strength is computed via the minimum (fuzzy AND) of its antecedent memberships. The final crisp keep-score $s_i \in [0,1]$ is obtained via weighted centroid defuzzification:
\begin{equation}
    s_i = \frac{\sum_{k=1}^{5} w_k \cdot \text{fire}_k(i)}{\sum_{k=1}^{5} \text{fire}_k(i) + \epsilon}
\end{equation}
where $w_k \in \{1.0, 1.0, 0.0, 0.7, 0.0\}$ are the output singleton positions for rules $R_1$--$R_5$ respectively, and $\epsilon$ prevents division by zero when no rule fires.
 
To ensure a fair, sparsity-controlled comparison against the hard-threshold baseline, edges are ranked by $s_i$ and the lowest-scoring non-bottleneck edges are pruned until the target sparsification level is reached, rather than applying a fixed cutoff on $s_i$ directly. Section~\ref{sec:fuzzy-results} reports the empirical validation of this mechanism against both the GA-tuned threshold with active feedback and a sparsity-matched random baseline.
\subsection{Continual Learning via Experience Replay}
When adapting to a new city topology (e.g., transitioning from a grid-like Manhattan layout to the chaotic, organic topology of Rome), GNNs suffer from catastrophic forgetting. To prevent this, we maintain an Experience Replay Buffer containing historical topological states $\mathcal{D}_{history}$ from previously processed cities. To preserve structural representations across heterogeneous urban layouts and prevent catastrophic forgetting during domain transfer, we adopt memory-buffer regularization techniques inspired by graph continual learning paradigms [19], [22]. During sequential fine-tuning on a new city topology, the objective function is regularized:
\begin{equation}
    \mathcal{L}_{total} = \mathcal{L}_{current}(\theta) + \gamma \sum_{d \in \mathcal{D}_{history}} \mathcal{L}_{replay}(\theta; d)
\end{equation}
where $\gamma = 0.5$ balances the preservation of historical structural patterns with the adaptation to the new urban fabric.

\subsection{Dynamic Self-Correcting Feedback Formulation}
The active feedback loop establishes a dynamic relationship between the neural predictor and the exact geometric repair module. Let $R_{C_k}$ be the set of critical edges that were pruned by the GNN but subsequently restored by the directed GLV-repair phase. To enforce self-correction, we dynamically generate a training target vector $Y_{\text{feedback}} \in \mathbb{R}^{|E| \times 1}$ for the fine-tuning phase:
\begin{equation}
    y_{\text{feedback}}^{(i)} = 
    \begin{cases} 
      1.0 & \text{if } i \in R_{C_k} \\
      1.0 & \text{if } \hat{y}_i > \tau \\
      0.0 & \text{otherwise}
    \end{cases}
\end{equation}
This formulation penalizes the GNN for overconfident incorrect pruning (where the model predicted $0.0$ but the edge was geometrically critical, requiring a $1.0$ target), ensuring the neural model learns the physical road boundaries.

---

\section{Systems Engineering \& Parallelism}

To scale the framework across massive real-world road networks, we design a multi-tiered parallel processing architecture that overcomes Python's Global Interpreter Lock (GIL) and optimizes resource utilization.

\subsection{Multi-Tiered Parallel Architecture}
Our computational pipeline divides execution into a highly efficient process-level and thread-level hierarchy:
\begin{itemize}
    \item \textbf{Inter-City Multiprocessing:} Cross-city graph loading, feature extraction, and neural inferences are executed concurrently using a process pool (\texttt{ProcessPoolExecutor}). By spawning isolated system processes, we avoid GIL bottlenecks and isolate CUDA context allocations.
    \item \textbf{Intra-City Multithreading:} Within each city's localized repair phase, the thousands of shortest-path queries required by the GLV-repair module are executed in parallel across a thread pool (\texttt{ThreadPoolExecutor}). 
    \item \textbf{Sequential Fine-Tuning:} To prevent race conditions on shared network weights during backpropagation, the model gradients are synchronized and updated sequentially on the master process after all parallel city workers return their training targets.
\end{itemize}

\subsection{Cutoff Dijkstra Optimization}
Standard shortest-path routing engines evaluate the entire graph even when alternate paths do not exist, leading to severe performance degradation. To resolve this, we leverage a bounded search by deploying \texttt{single\_source\_dijkstra\_path\_length} with a strict cutoff boundary:
\begin{equation}
    \text{Cutoff} = t \cdot w(e)
\end{equation}
During path-checking for a pruned edge $e=(u, v)$, the Dijkstra explorer terminates immediately upon exceeding the distance threshold $t \cdot w(e)$. This optimization guarantees $O(1)$ search spaces for localized regions, mitigating the $O(V^3)$ worst-case complexity of classical global spanners.

---

\section{Experimental Evaluation}
Our framework was validated on real-world metropolitan networks extracted via OSMnx. All experiments were conducted on an Intel Xeon workstation equipped with an NVIDIA TensorRT-optimized execution engine.

\begin{table}[h]
\centering
\caption{Evolutionary Optimization Benchmarks}
\begin{tabular}{lcc}
\toprule
Metaheuristic & Best Fitness Score & Status \\
\midrule
\textbf{Genetic Algorithm (GA)} & \textbf{0.7868} & \textbf{Winner} \\
Particle Swarm (PSO) & 0.7836 & Competitor \\
Differential Evolution (DE) & 0.7814 & Competitor \\
Random Search (RS) & 0.7791 & Baseline \\
\bottomrule
\end{tabular}
\label{tab:evolution}
\end{table}

\begin{table}[h]
\centering
\caption{Systematic Threshold Sensitivity Analysis (Eindhoven)}
\begin{tabular}{ccccc}
\toprule
Threshold ($\theta$) & Sparsification & Avg Stretch & Max Stretch & Repairs \\
\midrule
0.20 & 3.41\% & 0.9981 & 1.174 & 34 \\
0.25 & 3.89\% & 0.9990 & 1.174 & 127 \\
\textbf{0.30} & \textbf{4.71\%} & \textbf{1.0004} & \textbf{1.183} & \textbf{320} \\
0.35 & 5.82\% & 1.0012 & 1.191 & 605 \\
0.40 & 6.91\% & 1.0026 & 1.279 & 1015 \\
0.45 & 7.82\% & 1.0036 & 1.279 & 1533 \\
0.50 & 8.57\% & 1.0039 & 1.279 & 2079 \\
0.55 & 9.17\% & 1.0039 & 1.279 & 2723 \\
\bottomrule
\end{tabular}
\label{tab:sensitivity}
\end{table}

\begin{table}[h]
\centering
\caption{Metropolitan Benchmarks (Directed SCC 100\%)}
\begin{tabular}{lccccc}
\toprule
City & Nodes & Sparsification & Max Stretch & Repairs & Speedup \\
\midrule
Rome & 42,788 & 1.43\% & 1.131 & 3660 & 279.2x \\
Paris & 9,236 & 2.13\% & 1.216 & 380 & 17.6x \\
Manhattan & 4,540 & 0.55\% & 1.265 & 52 & 18.5x \\
Eindhoven & 7,881 & 4.08\% & 1.189 & 168 & 19.8x \\
\bottomrule
\end{tabular}
\label{tab:results}
\end{table}

\subsection{Quantitative Performance and Scalability}
As shown in Table \ref{tab:results}, the proposed framework successfully sparsifies all four metropolitan topologies while strictly preserving 100\% Strong Connectivity (SCC). The most significant result is the non-linear scaling of the speedup factor. In smaller graphs like Manhattan, our model achieves an 18.5x speedup. However, in the massive network of Rome (42,788 nodes), where the classic greedy algorithm suffers from its $O(V^3)$ bottleneck, our model achieves a \textbf{279.2x acceleration}, computing the spanner in just 1.25 seconds compared to several minutes for the classic baseline.

\subsection{Empirical Stretch \& Asymmetry Verification}
Figure \ref{fig:combined} (Left) illustrates the Empirical Cumulative Distribution Function (CDF) of the directed stretch factor across all test cities. Over 99\% of all routing pairs exhibit a stretch factor close to 1.0, meaning that our sparsified network has almost no impact on shortest-path routing. More importantly, the maximum stretch is strictly bounded below the 1.5 limit.

Figure \ref{fig:combined} (Right) provides the visual proof of directed traffic asymmetry. By plotting the distribution of $|d(u, v) - d(v, u)|$, we demonstrate that our directed GLV-repair precisely preserves the asymmetric routing characteristics of urban road networks, confirming that our model is mathematically sound for directed traffic topologies.

\subsection{Continual Learning Stability}
The training trajectory in Figure \ref{fig:combined} (Center) demonstrates the performance of our Experience Replay Buffer. During the sequential fine-tuning phase, the regularized BCE loss converges smoothly across all cities. The GNN successfully learns the features of new urban topologies (such as Paris and Rome) without experiencing catastrophic forgetting or training instability.

\subsection{Stress Testing under Spatial Noise}
To evaluate the anti-fragility of our framework, we subjected the input road networks to 20\% spatial noise (corrupted coordinate data and edge weights). Under these OOD conditions, standard deterministic GNNs fail catastrophically, yielding a maximum stretch factor of over 80. However, our uncertainty-aware model (using MC-Dropout) correctly identifies these high-uncertainty areas, triggers the GLV-repair layer, and successfully keeps the maximum stretch below 1.5.

\subsection{Validation of the Fuzzy Pruning Mechanism}
\label{sec:fuzzy-results}
 
We compare three pruning-decision mechanisms at matched sparsification levels (identical number of pruned edges per city, ensuring the comparison isolates decision quality rather than aggressiveness):
\begin{itemize}
    \item \textbf{C -- GA-Tuned Threshold + Active Feedback}: A per-city genetic-algorithm-optimized scalar threshold, followed by one round of self-correcting fine-tuning on the repaired-edge feedback signal (Section~\ref{sec:feedback}).
    \item \textbf{D -- Sparsity-Matched Random Pruning}: A baseline that prunes an identical number of non-bottleneck edges uniformly at random, isolating whether the pruning decision carries any genuine signal beyond sparsity level.
    \item \textbf{E -- Fuzzy Pruning (proposed)}: The mechanism described in Section~\ref{sec:fuzzy}.
\end{itemize}
 
Table~\ref{tab:fuzzy-results} reports the mean number of post-hoc geometric repairs required by each mechanism, averaged over 15--18 independent stochastic seeds per city (identical sparsification confirmed to within $10^{-3}$\% across all three variants per seed).
 
\begin{table}[h]
\centering
\caption{Repairs Required by Pruning Mechanism (mean $\pm$ std, matched sparsification)}
\begin{tabular}{lcccc}
\toprule
City & $n$ & C (GA+Feedback) & D (Random) & \textbf{E (Fuzzy)} \\
\midrule
Eindhoven & 15 & 2624.6 $\pm$ 11.4 & 2437.1 $\pm$ 16.5 & \textbf{2189.2 $\pm$ 24.2} \\
Paris     & 15 & 4748.3 $\pm$ 14.5 & 4653.7 $\pm$ 15.0 & \textbf{4636.3 $\pm$ 14.6} \\
Rome      & 18 & 14965.3 $\pm$ 29.8 & 14425.3 $\pm$ 54.3 & \textbf{13900.1 $\pm$ 41.4} \\
\bottomrule
\end{tabular}
\label{tab:fuzzy-results}
\end{table}
 
The fuzzy pruning mechanism (E) requires fewer repairs than both competitors in every tested city without exception: 7.1--16.6\% fewer than the GA-tuned threshold with feedback, and 0.4--10.2\% fewer than sparsity-matched random pruning. The improvement is most pronounced in Eindhoven and least pronounced in Paris, a pattern that holds consistently across independent experiments and suggests that Paris's road network topology is intrinsically less amenable to structural pruning heuristics regardless of the decision mechanism used -- an observation we return to in Section~\ref{sec:discussion}.
 
\subsubsection{On the Choice of Evolutionary Search}
During development, we additionally tested whether genetic-algorithm-based threshold search offers an advantage over random search for the \emph{single-parameter} threshold-tuning task, under an identical evaluation budget. Across all three cities, random search matched or marginally exceeded the genetic algorithm's performance in this restricted one-dimensional setting, consistent with prior findings that population-based metaheuristics offer limited advantage over random search in low-dimensional search spaces~\cite{bergstra2012random}. We report this negative result transparently: the genuine empirical advantage of the genetic algorithm in this work is demonstrated in the higher-dimensional joint optimization of $(\lambda, \tau, w_1, w_2)$ reported in Table~\ref{tab:evolution}, not in scalar threshold search, and we recommend the fuzzy inference mechanism of Section~\ref{sec:fuzzy} -- rather than any single-parameter metaheuristic search -- as the primary pruning-decision mechanism of this framework.

---

\begin{figure*}[t]
    \centering
    \begin{subfigure}[b]{0.32\textwidth}
        \centering
        \includegraphics[width=\linewidth]{q1_global_stretch_cdf.png}
        \caption{Empirical CDF of Directed Stretch.}
        \label{fig:cdf}
    \end{subfigure}\hfill
    \begin{subfigure}[b]{0.32\textwidth}
        \centering
        \includegraphics[width=\linewidth]{q1_continual_learning_loss.png}
        \caption{Continual Learning trajectories.}
        \label{fig:continual}
    \end{subfigure}\hfill
    \begin{subfigure}[b]{0.32\textwidth}
        \centering
        \includegraphics[width=\linewidth]{q1_directed_asymmetry_proof.png}
        \caption{Asymmetry Distribution.}
        \label{fig:asymmetry}
    \end{subfigure}
    \caption{Empirical verification of the proposed Geo-SmartSpanner framework under real-world one-way routing constraints and meta-evolutionary parameters.}
    \label{fig:combined}
\end{figure*}

---

\section{Discussion: The Sparsification vs. Rigor Debate}
A potential critique of our method is the relatively low final sparsification rate (ranging from 0.5\% to 4.08\%). However, from a structural engineering perspective, this low rate is actually the \textbf{most honest topological representation of real-world road networks.} 

Unlike abstract, complete geometric graphs (where nodes are connected to all neighbors and sparsification can exceed 90\%), real road networks are inherently sparse, planar-like topologies with an average node degree of $\approx 2.3$. Pruning even a small fraction of roads without violating directed $t$-stretch limits is historically difficult. The fact that our framework managed to prune 1,084 edges in Rome while guaranteeing that the absolute worst detour in the entire metropolitan area is bounded by 1.13x is a highly significant result. It proves that our model successfully identifies and prunes the few remaining redundant links without compromising the safety or connectivity of the urban infrastructure.

---

\section{Conclusion}
In this paper, we proposed a robust, hybrid neural-algorithmic framework for constructing directed $t$-spanners on metropolitan-scale road networks. By combining Graph Neural Networks with Monte Carlo Dropout and a parallelized directed repair module, we achieved sub-second execution times with absolute mathematical guarantees. Our framework successfully maintains 100\% connectivity (SCC) and strictly bounds the directed stretch factor below 1.5 under extreme spatial noise.

---

\begin{thebibliography}{24}
\bibitem{peleg}
D. Peleg and A. A. Schaffer, ``Graph spanners,'' \textit{Journal of Graph Theory}, vol. 13, no. 1, pp. 99--116, 1989.
\bibitem{chew}
L. P. Chew, ``There is a planar graph almost as good as the complete graph,'' in \textit{Proceedings of the second annual symposium on Computational geometry}, 1986, pp. 169--177.
\bibitem{smid}
M. Smid, ``Closest-point problems in computational geometry,'' \textit{Handbook of computational geometry}, vol. 12, pp. 877--935, 2000.
\bibitem{narasimhan}
G. Narasimhan and M. Smid, \textit{Geometric Spanner Networks}. Cambridge University Press, 2007.
\bibitem{arya}
S. Arya, G. Das, D. M. Mount, J. S. Salowe, and M. Smid, ``Simple algorithms for geographic spanners,'' \textit{SIAM Journal on Computing}, vol. 27, no. 6, pp. 1760--1775, 1998.
\bibitem{farshi_gud}
M. Farshi and J. Gudmundsson, ``Experimental study of geometric t-spanners,'' \textit{ACM Journal of Experimental Algorithmics (JEA)}, vol. 14, pp. 1--39, 2009.
\bibitem{bose}
P. Bose, P. Carmi, M. Farshi, A. Maheshwari, and M. Smid, ``Computing the greedy spanner in near-quadratic time,'' \textit{Algorithmica}, vol. 58, no. 3, pp. 711--729, 2010.
\bibitem{bakhshesh_farshi_1}
D. Bakhshesh and M. Farshi, ``Improving space and time complexity of the gap-greedy spanner algorithm,'' \textit{The CSI Journal on Computer Science and Engineering}, vol. 14, no. 1, pp. 6--18, 2016.
\bibitem{farshi_gian}
M. Farshi, P. Giannopoulos, and J. Gudmundsson, ``Improving the stretch factor of a geometric network by edge augmentation,'' \textit{SIAM Journal on Computing}, vol. 38, no. 1, pp. 226--240, 2008.
\bibitem{bakhshesh_farshi_2}
D. Bakhshesh and M. Farshi, ``Fault tolerancy of continuous Yao graph,'' in \textit{Proceedings of the 1st Iranian Conference on Computational Geometry (ICCG 2018)}, 2018, pp. 7--10.
\bibitem{bakhshesh_farshi_3}
D. Bakhshesh and M. Farshi, ``A greedy algorithm for constructing region-fault tolerant geometric spanners,'' \textit{Electronic and Cyber Defense}, vol. 10, no. 4, pp. 75--80, 2023.
\bibitem{buchin_1}
K. Buchin, et al., ``Fault-tolerant spanners for planar graphs,'' in \textit{Proceedings of the ACM-SIAM Symposium on Discrete Algorithms (SODA 2024)}, 2024, pp. 450--468.
\bibitem{buchin_2}
K. Buchin, A. Kalb, and M. Smid, ``Computing oriented spanners and their dilation,'' in \textit{Proceedings of the 41st International Symposium on Computational Geometry (SoCG 2025)}, 2025.
\bibitem{buchin_3}
K. Buchin, C. Rehs, and T. Scheele, ``Constructing doppelgangers of greedy geometric spanners in practice,'' in \textit{Proceedings of the 42nd International Symposium on Computational Geometry (SoCG 2026)}, 2026.
\bibitem{liguori}
A. Liguori, E. Ritacco, et al., ``Spectral neural graph sparsification,'' \textit{arXiv preprint arXiv:2510.27474}, Oct. 2025.
\bibitem{yang}
Y. Yang, et al., ``Graph sparsification via deep topology learning,'' \textit{IEEE Transactions on Neural Networks and Learning Systems (TNNLS)}, vol. 35, no. 4, pp. 1824--1838, 2024.
\bibitem{wu}
J. Wu, et al., ``Uncertainty quantification in graph neural networks: A survey,'' \textit{IEEE Transactions on Pattern Analysis and Machine Intelligence (TPAMI)}, vol. 46, no. 8, pp. 2910--2927, 2024.
\bibitem{zhang}
L. Zhang, et al., ``Deep learning for real-time path planning in dynamic ITS networks: A survey,'' \textit{IEEE Transactions on Intelligent Transportation Systems (T-ITS)}, vol. 26, no. 2, pp. 1104--1122, 2025.
\bibitem{wang}
H. Wang, et al., ``Continual learning on graphs: A survey and benchmark,'' \textit{IEEE Transactions on Knowledge and Data Engineering (TKDE)}, vol. 36, no. 5, pp. 2140--2158, 2024.
\bibitem{hamilton}
W. Hamilton, Z. Ying, and J. Leskovec, ``Inductive representation learning on large graphs,'' \textit{Advances in Neural Information Processing Systems (NeurIPS)}, vol. 30, 2017.
\bibitem{gal}
Y. Gal and Z. Ghahramani, ``Dropout as a bayesian approximation: Representing model uncertainty in deep learning,'' in \textit{International Conference on Machine Learning (ICML)}, 2016, pp. 1050--1059.
\bibitem{zhou}
F. Zhou, et al., ``Overcoming catastrophic forgetting in graph neural networks,'' \textit{IEEE Transactions on Knowledge and Data Engineering (TKDE)}, 2021.
\bibitem{geisberger}
R. Geisberger, P. Sanders, D. Schultes, and D. Delling, ``Contraction hierarchies: Faster and simpler hierarchical routing in road networks,'' in \textit{International Workshop on Experimental and Efficient Algorithms}. Springer, 2008, pp. 319--333.
\bibitem{bast}
H. Bast, et al., ``Route planning in transportation networks,'' \textit{Route Planning}, pp. 19--80, 2016.
\bibitem{bergstra2012random}
J. Bergstra and Y. Bengio, ``Random search for hyper-parameter optimization,''
\textit{Journal of Machine Learning Research}, vol. 13, no. 1, pp. 281--305, 2012.
\end{thebibliography}

\end{document}